\clearpage
\fancyhead[L]{\small DE LATITUDINE SPECULAE MANHEMIENSIS.}
\hypertarget{native-602-lower}{}
\begin{center}\small Göttingische gelehrte Anzeigen. Stück 76. Seite 753..756. 1817 Mai 12.\\[3pt]\rule{.72\linewidth}{.3pt}\end{center}

\emph{De latitudine speculae Manhemiensis autore} H.\,C.\ \textsc{Schumacher},
\emph{astronomiae professore et regiae societatis scientiarum Havniensis
socio.} Copenhagen 1816. Gedruckt bei J. F. Schulz. 56 Seiten in gr. Quart.

Die Sternwarte in Mannheim besitzt, zur Messung der Zenithdistanzen, einen
achtfussigen Mauer\SourceLeaf{603}quadranten von \textsc{Bird}, einen neunfussigen Zenithsector von \textsc{Sisson}, und seit dem Jahr 1811 auch einen
dreifussigen Vervielfältigungskreis von \textsc{Reichenbach}. Zur Aufstellung des letztern Instruments sind aber
bisher noch keine Veranstaltungen getroffen, und der Verf., von 1813 bis 1815 Director jener Sternwarte,
war daher, um die Polhöhe derselben so genau als es sich mit den vorhandenen Hülfsmitteln thun liess,
zu bestimmen, auf die beiden ersten Instrumente beschränkt. Der Zenithsector diente zur Ausmittlung
des Collimationsfehlers des Mauerquadranten; allein da dieser nach Süden gerichtet ist und nicht umgehängt werden kann, so war eine selbstständige Bestimmung der Polhöhe durch Fixsterne unmöglich,
und es blieb nichts übrig, als die Declinationen der beobachteten Fixsterne nach den zuverlässigsten
und neuesten Bestimmungen anderer Astronomen zum Grunde zu legen. Einhundert und zwei und funfzig am Zenithsector vom 27. Januar 1814 bis 4. März 1815 angestellte Beobachtungen dienten dazu, von
achtzehn Sternen die mittleren auf den Anfang des Jahrs 1815 reducirten Zenithdistanzen zu bestimmen.
Diese bedurften jedoch einer Correction, da sich nach angestellter Prüfung ergab, dass die Grade auf
dem Zenithsector zu klein, mithin die gemessenen Zenithdistanzen zu gross waren. Um das Gesetz dieser Correctionen auszumitteln, mass der Verf. mit einem Stangenzirkel die Chorden von mehrern Bögen auf dem Limbus des Sectors und zugleich die Entfernung von dem Mittelpunkte der Bewegung.
Die aus diesen Datis von Hrn. \textsc{Bessel} in Königsberg abgeleitete Formel für die Correction
\[
1''3892\,z
\]
wo $z$ die Zenithdistanz in Graden bedeutet, stellt jene Messungen ziemlich gut dar, und die kleinen
Differenzen können füglich als Fehler der Messungen angesehen werden. Der Verf. hat inzwischen die
Formel
\[
0''4225\left(z+\tfrac{1}{10}zz+\tfrac{1}{20}z^{3}\right)
\]
vorgezogen, die eine noch bessere Uebereinstimmung zeigt. Hiebei ist also angenommen, dass die Grade
des Zenithsectors nicht blos zu klein, sondern auch von ungleicher Grösse sind; genau genommen ist
jedoch das zweite Glied aus dem Grunde hier nicht zulässig, weil alle die Bögen, deren Chorden gemessen wurden, vom Nullpunkte halbirt werden; auch würde die Formel
\[
0''585\left(z+\tfrac{1}{28}z^{3}\right)
\]
eine eben so gute Uebereinstimmung gegeben haben: übrigens ist es für das Endresultat fast ganz gleichgültig, welcher von diesen Correctionsformeln man sich bediene. Die so gefundenen Zenithdistanzen
dienten nun zur Bestimmung des Collimationsfehlers des Mauerquadranten, an welchem dieselben Sterne
beobachtet worden waren, und hiermit wurden die Zenithdistanzen von 36 andern am Mauerquadranten
beobachteten Sterne berechnet, deren Vergleichung mit den von \textsc{Piazzi}, \textsc{Pond} und \textsc{Oriani} bestimmten
Declinationen für die Polhöhe der Mannheimer Sternwarte, als Resultat aus 284 Beobachtungen $49^{\circ}\,29'\,13''5$
gaben. Von den 18 am Zenithsector beobachteten Sternen selbst konnte der Verf. nur neun auf gleiche
Weise benutzen, da ihm gleich zuverlässige Declinationsbestimmungen der übrigen fehlten. Jene gaben
ihm aus 87 Beobachtungen die Polhöhe $49^{\circ}\,29'\,14''2$. Recens. hat die Beobachtungen der übrigen neun,
nach den Declinationsbestimmungen in \textsc{Piazzi}'s neuem Catalog (welchen der Verf. noch nicht benutzen
konnte) berechnet, und folgende Resultate gefunden:
\begin{center}
\fontsize{9}{11}\selectfont
\setlength{\tabcolsep}{3pt}
\begin{tabular}{l|c|c||l|c|c}
 & Polhöhe & \begin{tabular}{c}Anzahl der\\Beobachtungen\end{tabular} & & Polhöhe & \begin{tabular}{c}Anzahl der\\Beobachtungen\end{tabular} \\
\hline
$\sigma$ Perseus & $49^{\circ}\,29'\,10''4$ & 8 & 725 Perseus & $49^{\circ}\,29'\,11''5$ & 9 \\
$\psi$ Perseus & 8.3 & 9 & 7 Perseus & 8.2 & 8 \\
$\chi$ gr. Bär & 8 3 & 5 & $F$ gr. Bär & 9.8 & 5 \\
$\theta$ Perseus & 15.2 & 5 & $\gamma$ Perseus & 8.2 & 13 \\
$\iota$ gr. Bär & 8.0 & 5 & \multicolumn{3}{|l}{} \\
\end{tabular}
\end{center}
\noindent\SourceLeaf{604}Für den von dem Verf. als 725 Perseus bezeichneten Stern, dessen Designation nicht aufzufinden war, wurde in \textsc{Piazzi}'s Catalog II, 253 angenommen. Diese 65 Beobachtungen geben also im Mittel $49^\circ\,29'\,9''8$; folglich alle 152 Beobachtungen am Zenithsector im Mittel $49^\circ\,29'\,12''3$, also die sämmtlichen 436 Fixsternbeobachtungen im Mittel $49^\circ\,29'\,13''1$. Die von dem Verf. im Jahre 1814 angestellten Beobachtungen der Sonne, 66 an der Zahl, geben im Mittel nur einige Zehntel der Secunde mehr; sie können jedoch wegen einiger localen Ursachen nicht auf gleiche Zuverlässigkeit Anspruch machen. Dies Resultat aus den Beobachtungen des Verf. wird nun für das sicherste gelten müssen, bis der \textsc{Reichenbach}'sche Kreis aufgestellt und zu der einer solchen Sternwarte eigentlich allein würdigen neuen Bestimmung der Polhöhe durch Circumpolarsterne benutzt sein wird  Uebrigens kommt obiges Resultat zwar genau mit \textsc{Barry}'s letzter Angabe überein; allein da dieser Astronom die Art, wie er diese Bestimmung erhalten hat, nicht bekannt gemacht hatte, und man daher über den Grad der Genauigkeit derselben kein Urtheil fällen konnte, so verdient Hr. \textsc{Schumacher} für die vollständige Bekanntmachung seiner Beobachtungen den Dank aller Astronomen.
